\section{生徒の構成}
\label{sec_student}
本章では，教師ラベルで学習する生徒の構成と，学習データの集め方および損失を検討した結果を述べる．

\subsection{生徒のモデル}
生徒には，埋め込みに対する MLP と，微調整した画像エンコーダの2種類を用いた．
MLP の入力は，SigLIP2-so400m の画像埋め込みを L2 正規化した1,152次元のベクトル $ {\mathbf x} \in {\mathbb R}^{1 \times 1152} $ である．
MLP の中間の表現 $ {\mathbf h} $ と，各軸の出力を(\ref{eq_mlp})に示す．

\begin{align}
    {\mathbf h} &= {\rm Dropout} \left( {\rm GELU} \left( {\rm LN}({\mathbf x}) {\mathbf W}_1 + {\mathbf b}_1 \right) \right),
    \nonumber \\
    {\mathbf p}^{\rm F} &= {\rm softmax} \left( {\mathbf h} {\mathbf W}_{\rm F} + {\mathbf b}_{\rm F} \right),
    \nonumber \\
    {\mathbf p}^{\rm A} &= \sigma \left( {\mathbf h} {\mathbf W}_{\rm A} + {\mathbf b}_{\rm A} \right),
    \quad
    {\mathbf p}^{\rm M} = \sigma \left( {\mathbf h} {\mathbf W}_{\rm M} + {\mathbf b}_{\rm M} \right)
    \label{eq_mlp}
\end{align}

\noindent
ここで ${\rm LN}$ は層正規化\cite{ref_layernorm}，${\rm GELU}$ は活性化関数\cite{ref_gelu}，${\rm Dropout}$ は確率0.3のドロップアウト\cite{ref_dropout}，$ \sigma $ は要素ごとのシグモイド関数である．
$ {\mathbf W}_1 \in {\mathbb R}^{1152 \times 1024} $ と $ {\mathbf b}_1 \in {\mathbb R}^{1 \times 1024} $ が中間層の重み行列としきい値ベクトルである．
$ {\mathbf W}_{\rm F}, {\mathbf W}_{\rm A}, {\mathbf W}_{\rm M} $ と $ {\mathbf b}_{\rm F}, {\mathbf b}_{\rm A}, {\mathbf b}_{\rm M} $ は，形式，属性，媒体の出力層の重み行列としきい値ベクトルである．
$ {\mathbf p}^{\rm F} \in {\mathbb R}^{1 \times 4} $ は形式の4クラスの確率，$ {\mathbf p}^{\rm A} \in {\mathbb R}^{1 \times 3} $ と $ {\mathbf p}^{\rm M} \in {\mathbb R}^{1 \times 7} $ は属性と媒体の各項目が真である確率である．
微調整のモデルは，SigLIP2-B の画像エンコーダの出力に，形式，属性，媒体の3つの線形の出力層を付けたものである．

学習の損失 $ L $ は，形式の交差エントロピーと，属性と媒体の各項目の2値交差エントロピーの和とし，(\ref{eq_loss})に示す．

\begin{align}
    L = &- \sum_{k=1}^{4} y^{\rm F}_k \log p^{\rm F}_k
    \nonumber \\
    &- \sum_{j \in {\mathbb J}} \left\{ y_j \log p_j + (1 - y_j) \log (1 - p_j) \right\}
    \label{eq_loss}
\end{align}

\noindent
ここで $ y^{\rm F}_k \in \{0, 1\} $ は教師ラベルの形式がクラス $ k $ であるかを表し，$ {\mathbb J} $ は属性と媒体の項目の集合，$ y_j \in \{0, 1\} $ と $ p_j $ は項目 $ j $ の教師ラベルと予測の確率である．
最適化には AdamW\cite{ref_adamw} を用いた．
MLP は学習率 $ 3 \times 10^{-4} $ で40エポック学習し，重み減衰は $ 10^{-4} $ と $ 10^{-2} $ から教師ラベルの検証用分割で選んだ．
微調整は，入力を $ 224 \times 224 $ 画素とし，学習率 $ 3 \times 10^{-5} $，重み減衰0.05，バッチサイズ128で5エポック学習した．
文字の向きを保つため左右反転の拡張は用いず，色の揺らぎのみを拡張とした．
いずれも検証用分割での3軸の平均が最良のエポックの重みを用いた．

予測では，形式は確率が最大のクラスとし，属性と媒体は確率0.5以上の項目を真とした．
媒体が1つも真でない場合は，確率が最大の媒体を真とした．
最後に，教師と同じ規則で形式から属性を導いた．

\subsection{少ない区分の追加}
無作為抽出の教師ラベル40,000枚では，CG が305枚，絵画が217枚，カルタが494枚しか含まれず，生徒はこれらをほとんど学習できなかった．
\ref{sec_selection}章では，少ない種類を多めに含む画像に教師ラベルを付ける方が効率が良いことを示した．
そこで，次の手順で少ない区分の画像を追加した．
まず，学習にも評価にも用いていない30万枚の候補に SigLIP2-so400m の埋め込みを計算した．
次に，40,000枚で学習した MLP で候補ごとに各区分の確率を求め，目標に達していない区分について確率の高い順に約1,000枚を選び，教師ラベルを付けた．
教師がその区分と判定した画像が40,000枚との合計で $ N_{\rm min} = 1{,}000 $ 枚以上になるまで，これを繰り返した．
以下では，40,000枚のみの学習データを基本の学習データ，追加分を含む学習データを追加ありの学習データと呼ぶ．

5回の繰り返しで計4,998枚を追加し，全区分が1,000枚以上になった．
最終的な枚数は，ゲーム1,715枚，CG 1,161枚，絵画1,075枚，カルタ1,178枚であった．
絵画は最も集まりにくく，2,666枚に教師ラベルを付ける必要があった．
追加のラベル付けは約30分で済んだ．

\subsection{重み付きの損失}
少ない区分の学習を強めるため，クラスの出現率に応じて重みを付けた損失も比較した．
形式のクラス $ k $ の交差エントロピーの重み $ v_k $ と，属性と媒体の項目 $ j $ の正例の重み $ u_j $ を(\ref{eq_weight})に示す．

\begin{equation}
    v_k \propto \left( \frac{N}{N_k} \right)^{\alpha},
    \qquad
    u_j = \left( \frac{N - P_j}{P_j} \right)^{\alpha}
    \label{eq_weight}
\end{equation}

\noindent
ここで $ N $ は学習データの枚数，$ N_k $ は教師ラベルの形式がクラス $ k $ である画像の数，$ P_j $ は項目 $ j $ の正例の数，$ \alpha = 0.5 $ である．
$ v_k $ は平均が1になるように正規化した．

表\ref{tab_rare_weight}に，第2版の定義における微調整した生徒の結果を示す．
少ない区分の F1 は，少ない区分の評価用データ第1版における4区分の F1 の平均である．

\begin{table}[t]
    \centering
    \caption{少ない区分の追加と重み付きの損失の効果（第2版の定義，微調整した SigLIP2-B）}
    \label{tab_rare_weight}
    \small
    \begin{tabular}{lccc}
        \toprule
        条件 & 少ない区分 & 3軸の平均 & 頑健な指標 \\
        \midrule
        基本，重みなし & 0.528 & 0.781 & 0.857 \\
        基本，重みあり & 0.676 & 0.803 & 0.847 \\
        追加あり，重みなし & 0.632 & 0.803 & 0.854 \\
        追加あり，重みあり & 0.713 & 0.812 & 0.845 \\
        \midrule
        教師 & 0.699 & 0.810 & 0.848 \\
        \bottomrule
    \end{tabular}
\end{table}

表\ref{tab_rare_weight}の3軸の平均と頑健な指標は，通常の評価用データの test での値である．
少ない区分の追加と重み付きの損失は，いずれも単独で少ない区分の F1 を上げ，併用すると教師と同等の 0.713 に達した．
改善は主に CG と絵画の再現率によるものであり，CG の F1 は 0.15 から 0.44，絵画は 0.08 から 0.51 に向上した．
一方，頑健な指標は 0.857 から 0.845 にわずかに下がり，重み付けは少ない区分の再現率を上げる代わりに多い区分の適合率を下げた．

ところが，第3版の定義で同じ比較を行うと，重み付きの損失は効果がなかった．
test の媒体のマクロ F1 は，重みを付けると全ての条件で下がり，例えば追加ありの学習データで 0.896 から 0.840 に低下した．
また，媒体を2つ以上選ぶ画像の割合が，実データで6\%から7\%程度から8\%から11\%に増え，目視ラベルの6.3\%から離れた．
第2版で重みが有効であったのは，境界の曖昧な CG のような少ない区分があったためと考えられる．
そこで第3版以降では重み付きの損失を廃止し，少ない区分の追加のみを用いた．

\subsection{生徒のモデルの比較と確率の平均}
最終の教師ラベルで，生徒のモデルとその組み合わせを比較した．
組み合わせでは，複数のモデルの確率を(\ref{eq_ensemble})のように平均した．

\begin{equation}
    \bar{\mathbf p} = \frac{1}{M} \sum_{m=1}^{M} {\mathbf p}^{(m)}
    \label{eq_ensemble}
\end{equation}

\noindent
ここで $ M $ は組み合わせるモデルの数，$ {\mathbf p}^{(m)} $ は $ m $ 番目のモデルが出力した形式，属性，媒体の確率を連結したベクトルである．
予測の規則は，平均した確率 $ \bar{\mathbf p} $ に適用する．

選択の指標は，通常の評価用データの dev の3軸の平均と，少ない区分の評価用データの dev の4区分の F1 の平均の平均とした．
test と実データの目視ラベルは確認のみに用いた．
表\ref{tab_student_models}に主な結果を示す．
単体のモデルは3つの乱数の種の平均であり，「MLP 3種」は3つの乱数の種で学習した MLP の確率の平均である．

\begin{table*}[t]
    \centering
    \caption{生徒のモデルの比較（最終の教師ラベル）}
    \label{tab_student_models}
    \footnotesize
    \begin{tabular}{lcccccc}
        \toprule
        モデル & 選択の指標 & dev 3軸 & 少ない区分 test & test 3軸 & 全項目一致 & 枚/秒 \\
        \midrule
        MLP 3種 ＋ 微調整 SigLIP2-B 224画素 & {\bf 0.872} & 0.919 & 0.888 & 0.919 & 0.838 & 約150 \\
        MLP ＋ 微調整 SigLIP2-B 224画素 ＋ 384画素 & 0.863 & 0.925 & 0.862 & 0.914 & 0.838 & 約115 \\
        MLP ＋ 微調整 SigLIP2-B 224画素 & 0.860 & 0.915 & 0.889 & 0.925 & 0.846 & 約150 \\
        MLP ＋ 微調整 SigLIP2-B 384画素 & 0.858 & 0.924 & 0.856 & 0.920 & 0.837 & 約135 \\
        MLP（単体） & 0.847 & 0.900 & 0.845 & 0.889 & 0.797 & 約180 \\
        微調整 SigLIP2-B 224画素（単体） & 0.836 & 0.905 & 0.839 & 0.896 & 0.795 & 約750 \\
        微調整 SigLIP2-B 384画素（単体） & 0.825 & 0.911 & 0.824 & 0.896 & 0.806 & 約410 \\
        \bottomrule
    \end{tabular}
\end{table*}

最も効果が大きかったのは，MLP と微調整のモデルの確率の平均であった．
実データの全項目一致の推定値は，単体の約0.80から約0.84に，少ない区分の test は0.84から0.89に向上した．
MLP は少ない区分に強く，微調整のモデルは通常の画像の形式と文字に強いため，誤りの傾向が異なる2つのモデルの平均が互いの弱点を補ったと考えられる．
入力を384画素にした微調整は，単体では dev が高いものの少ない区分の F1 が低く，組み合わせても224画素を上回らなかった．
MLP は埋め込みを共有するため，3つの乱数の種を平均しても計算量はほとんど増えない．
以上から，選択の指標が最良であった「MLP 3種と224画素の微調整 SigLIP2-B の確率の平均」を最終の生徒とした．
